% Figure : RAG calibre par fiche - deux collections
\begin{figure}[H]
\centering
\begin{tikzpicture}[node distance=10mm and 12mm, font=\small]

  % Collection 1
  \node[storage, minimum height=2.4cm, minimum width=3cm] (fiches) {};
  \node[font=\small\bfseries, above=1pt of fiches] {fiches\_poste};
  \node[font=\footnotesize, align=center] at (fiches.center)
       {fiche\_id\\\textit{(SHA1 12c)}\\texte\\embedding};

  % Collection 2
  \node[storage, right=24mm of fiches, minimum height=2.4cm, minimum width=3.4cm] (cands) {};
  \node[font=\small\bfseries, above=1pt of cands] {candidats\_evalues};
  \node[font=\footnotesize, align=center] at (cands.center)
       {candidat\_id\\profil\_brut\\score, source\\\textbf{fiche\_id}};

  % flèche relation
  \draw[flow, dashed] (fiches.east) -- node[above, font=\footnotesize\itshape]{lien fiche\_id} (cands.west);

  % procedure de lecture
  \node[noteblock, below=10mm of fiches.south west, anchor=north west, text width=12.5cm]
       {\textbf{Procédure d'interrogation} (A4 lors du scoring d'un nouveau candidat) :
        \begin{enumerate}\itemsep=1pt\topsep=1pt
        \item \textit{embed} la fiche courante ;
        \item interroger \texttt{fiches\_poste}, retenir les fiches dont la similarité cosinus \(\geq 0{,}5\) ;
        \item si aucune n'atteint le seuil \(\rightarrow\) retour \texttt{[]} (silence explicite) ;
        \item sinon, interroger \texttt{candidats\_evalues} avec \texttt{where=\{fiche\_id : \$in : [...] \}}.
        \end{enumerate}};
\end{tikzpicture}
\caption[Mémoire RAG calibrée par fiche]{Architecture de la mémoire RAG calibrée par fiche de poste. Les deux collections sont reliées par un identifiant stable \texttt{fiche\_id} qui filtre la recherche de candidats sur les fiches sémantiquement proches, écartant le biais inter-fiches du RAG naïf.}
\label{fig:rag_calibre}
\end{figure}
